\section{从离散符号到分布式表示}

\subsection{传统词义表示的困境}

在语言学中，词义通常被理解为\textbf{指称语义}（denotational semantics）：一个符号（signifier）与一个概念或事物（signified）之间的配对关系。例如，``tree''一词的含义就是世界上所有的树。

\begin{figure}[H]
\centering
\includegraphics[width=0.95\textwidth]{slides-images/slide-017.jpg}
\caption{指称语义：词义被定义为符号与概念/事物之间的映射关系\protect\footnotemark}
\end{figure}
\footnotetext{Slides 第18页。}

然而，在计算机中表示词义，传统方法面临严重困难。

\subsubsection{WordNet 的局限}

在前神经网络时代，最常用的词义资源是 \textbf{WordNet}——一种手工构建的词汇关系网络，记录了同义词、上下位关系等。

\begin{figure}[H]
\centering
\includegraphics[width=0.95\textwidth]{slides-images/slide-018.jpg}
\caption{WordNet 示例：展示了 ``good'' 的多个义项和 ``panda'' 的上位关系链\protect\footnotemark}
\end{figure}
\footnotetext{Slides 第19页。}

\begin{warningbox}{WordNet 的三大问题}
\begin{itemize}
    \item \textbf{缺乏细微差别}：WordNet 认为 ``proficient'' 是 ``good'' 的同义词，但``That was a good shot'' 不能替换为 ``That was a proficient shot''
    \item \textbf{严重不完整}：缺少现代词汇和俚语（如 wicked、badass、nifty 等）
    \item \textbf{人工维护成本高}：需要持续的人力投入来更新和扩展
\end{itemize}
\end{warningbox}

\subsubsection{One-hot 向量的问题}

在传统计算系统中，词被表示为 \textbf{one-hot 向量}：每个词在词表中有一个唯一位置，对应向量中只有该位置为 1，其余为 0。

\begin{figure}[H]
\centering
\includegraphics[width=0.95\textwidth]{slides-images/slide-020.jpg}
\caption{One-hot 向量的问题：motel 和 hotel 虽然语义相近，但它们的 one-hot 向量正交，点积为零\protect\footnotemark}
\end{figure}
\footnotetext{Slides 第21页。}

\begin{importantbox}{One-hot 表示的根本缺陷}
One-hot 向量是\textbf{局部表示}（localist representation）：每个词仅在向量的一个位置上被表示。这意味着：
\begin{itemize}
    \item 任意两个不同词的 one-hot 向量都是\textbf{正交}的（点积 = 0）
    \item 向量中不包含任何关于词义\textbf{相似性}的信息
    \item 无法从表示本身判断 ``motel'' 和 ``hotel'' 比 ``motel'' 和 ``chair'' 更相似
\end{itemize}
解决方案：\textbf{学习将相似性编码到向量本身中}。
\end{importantbox}

\subsection{分布式语义与词向量}

\subsubsection{分布式假说}

一种完全不同的语义观为这个问题提供了解决思路——\textbf{分布式语义}（distributional semantics）：

\begin{knowledgebox}{分布式假说的哲学渊源}
\textbf{J.R. Firth}（英国语言学家，1957）：``You shall know a word by the company it keeps.''（从一个词的上下文可以了解这个词。）\\[3pt]
这一思想可追溯到 \textbf{Wittgenstein} 的语言哲学：词的含义在于其\textbf{使用}（use），而非抽象的指称。
\end{knowledgebox}

\begin{figure}[H]
\centering
\includegraphics[width=0.95\textwidth]{slides-images/slide-021.jpg}
\caption{分布式语义示例：``banking'' 一词的含义可以从它频繁出现的上下文中推断\protect\footnotemark}
\end{figure}
\footnotetext{Slides 第22页。}

分布式语义的核心思想是：一个词的含义由它经常出现的\textbf{上下文}（surrounding words）决定。如果两个词经常出现在相似的上下文中，它们就有相似的含义。

\subsubsection{词向量（Word Vectors）}

基于分布式语义的思想，我们可以为每个词学习一个\textbf{稠密向量}（dense vector），使得语义相近的词具有相似的向量表示。

\begin{figure}[H]
\centering
\includegraphics[width=0.95\textwidth]{slides-images/slide-022.jpg}
\caption{词向量示例：banking 和 monetary 的 8 维稠密向量表示，相似的含义反映在相似的向量值上\protect\footnotemark}
\end{figure}
\footnotetext{Slides 第23页。}

\begin{importantbox}{词向量的核心特征}
\begin{itemize}
    \item \textbf{稠密}：每个维度都有非零值（区别于 one-hot 的稀疏表示）
    \item \textbf{低维}：典型维度为 100--2000（远小于词表大小，如 500,000）
    \item \textbf{分布式}：词义分布在整个向量的所有维度上（区别于 one-hot 的局部表示）
    \item \textbf{相似性可计算}：语义相似的词具有较大的\textbf{点积}（dot product）
\end{itemize}
词向量也被称为 \textbf{word embeddings}（词嵌入）或 \textbf{neural word representations}（神经词表示）。
\end{importantbox}

\subsubsection{词向量空间的可视化}

将词向量投影到二维空间（通常使用 t-SNE），可以看到语义相近的词聚集在一起：

\begin{figure}[H]
\centering
\includegraphics[width=0.95\textwidth]{slides-images/slide-024.jpg}
\caption{词向量空间的二维 t-SNE 可视化：国家名称聚集在一起，国籍词聚集在一起，动词按语义分组\protect\footnotemark}
\end{figure}
\footnotetext{Slides 第25页。}

\begin{knowledgebox}{高维空间的``神奇''特性}
Manning 特别强调：高维空间的行为与低维空间\textbf{截然不同}。在二维空间中，两个点只有当 x 和 y 坐标都相似时才``接近''。但在高维空间中，一个点可以\textbf{同时}在不同维度上靠近不同的点。

这就是为什么即使每个词只有\textbf{一个}向量，``star'' 也能同时靠近天文词汇（nebula、galaxy）和娱乐词汇（celebrity、fame）——这些相似性体现在不同的维度组合上。
\end{knowledgebox}

\begin{warningbox}{一词一向量的局限}
在 Word2Vec 中，每个词串（string）只有\textbf{一个}向量。这意味着多义词（如 ``bank''：银行/河岸，``star''：恒星/明星）的不同含义被\textbf{平均}到一个向量中。后续课程将介绍\textbf{上下文化词表示}（contextual word representations），如 ELMo、BERT 等，可以根据上下文给出不同的向量。
\end{warningbox}

\subsection{本章小结}

从 WordNet 的手工规则到 one-hot 的符号表示，传统方法都无法有效捕捉词义的细微差别和相似性。分布式语义假说提供了一条全新的路径：通过词的上下文来学习其含义。词向量将每个词表示为稠密的低维向量，使得语义相似性可以直接通过向量运算来度量。

%% ========== Part 4: Word2Vec ==========

